\documentclass[12pt]{article}
\usepackage[utf8]{inputenc}
\usepackage{amsmath, amssymb, amsthm}
\usepackage{logicproof}
\usepackage[most]{tcolorbox}
\usepackage{xcolor}
\usepackage{geometry}

\geometry{a4paper, margin=0.8in}

% Custom Colors (Teal Blue Flare)
\definecolor{tealblue}{HTML}{0EA5A4}
\definecolor{lightteal}{HTML}{E6F6F6}

\tcbset{
    colback=lightteal,
    colframe=tealblue,
    fonttitle=\bfseries,
    arc=3pt
}

\title{\color{tealblue} Homework 3}
\author{\textbf{Nguyen Dinh Thien Loc} \\ Student ID: 24125093}


\begin{document}

\maketitle

\section*{{\color{tealblue} 2.4.2} Model Satisfaction}
\textbf{Sentence:} $\phi \triangleq \forall x \exists y \exists z (P(x,y) \wedge P(z,y) \wedge (P(x,z) \rightarrow P(z,x)))$

\begin{tcolorbox}[title=Evaluation of Natural Number Models]
\begin{itemize}
    \item \textbf{(a) $\mathcal{M}: (\mathbb{N}, <)$}: 
    Let $x$ be arbitrary. Choose $y = x+1$ and $z = x$. 
    $P(x,y)$ is true as $x < x+1$. $P(z,y)$ is true as $x < x+1$. 
    For the implication, $P(x,z)$ is $x < x$ (False), making $F \rightarrow F$ true. 
    Thus, $\mathcal{M} \models \phi$.

    \item \textbf{(b) $\mathcal{M}': (\mathbb{N}, P(m,n) \text{ iff } n=2m)$}: 
    For any $x$, choose $y = 2x$ and $z = x$. 
    This satisfies $P(x,y)$ and $P(z,y)$. Since $x=z$, the final conjunct holds vacuously. 
    Thus, $\mathcal{M}' \models \phi$.

    \item \textbf{(c) $\mathcal{M}'': (\mathbb{N}, m < n+1)$}: 
    Choosing $y = x+1$ and $z = x$ satisfies $P(x,y)$ ($x < x+2$) and $P(z,y)$ ($x < x+2$). 
    The implication holds as $P(x,x)$ is $x < x+1$, but $P(z,x)$ is also $x < x+1$ (True $\rightarrow$ True).
    Thus, $\mathcal{M}'' \models \phi$.
\end{itemize}
\end{tcolorbox}

\section*{{\color{tealblue} 2.4.3} Identity and Predicates}
\textbf{Goal:} Find a model satisfying $\forall x \neg P(x,x)$ and one that does not.

\begin{itemize}
    \item \textbf{Model $\mathcal{M} \models \forall x \neg P(x,x)$}: Let $D=\{1,2\}$ and $P^{\mathcal{M}}=\{(1,2),(2,1)\}$. No element relates to itself.
    \item \textbf{Model $\mathcal{M}' \not\models \forall x \neg P(x,x)$}: Let $D=\{1\}$ and $P^{\mathcal{M}'}=\{(1,1)\}$. Here $P(1,1)$ is true, so the negation is false.
\end{itemize}

\section*{{\color{tealblue} 2.4.7} Semantic Entailment Proof}
\textbf{Show:} $\forall x \neg \phi \models \neg \exists x \phi$

\begin{proof}
    Assume $\mathcal{M} \models \forall x \neg \phi$. By definition, for every $a \in D^{\mathcal{M}}$, $\mathcal{M} \models \neg \phi[a/x]$, meaning $\phi$ is false for all $a$. 
    Suppose $\mathcal{M} \models \exists x \phi$. This requires at least one $a$ where $\phi[a/x]$ is true. 
    This is a direct contradiction. Thus, the model must satisfy $\neg \exists x \phi$.
\end{proof}

\section*{{\color{tealblue} 2.4.8} Semantic Entailment}
\textbf{Show:} $\forall x P(x) \vee \forall x Q(x) \models \forall x (P(x) \vee Q(x))$

\begin{proof}
    Assume $\mathcal{M} \models \forall x P(x) \vee \forall x Q(x)$.
    If $\mathcal{M} \models \forall x P(x)$, then for any $a$, $P(a)$ is true, making $P(a) \vee Q(a)$ true.
    If $\mathcal{M} \models \forall x Q(x)$, then for any $a$, $Q(a)$ is true, making $P(a) \vee Q(a)$ true.
    In both cases, $\forall x (P(x) \vee Q(x))$ holds.
\end{proof}

\section*{{\color{tealblue} 2.4.10} Counter-justification}
\textbf{Is $\forall x (P(x) \vee Q(x)) \models \forall x P(x) \vee \forall x Q(x)$? No.}

\begin{tcolorbox}
    Let $D = \{a, b\}$ with $P^{\mathcal{M}} = \{a\}$ and $Q^{\mathcal{M}} = \{b\}$.
    Every element in $D$ satisfies either $P$ or $Q$, so the left side is true.
    However, $\forall x P(x)$ is false (fails for $b$) and $\forall x Q(x)$ is false (fails for $a$).
    Thus, $T \models F$ is false.
\end{tcolorbox}

\section*{{\color{tealblue} 2.4.12} Validity and Counter-examples}

\begin{tcolorbox}[title=(a) Symmetry vs. Reflexivity]
\textbf{Invalid.} Let $D = \{1\}, S = \{(1,1)\}$. Premise $S(1,1) \to S(1,1)$ is True, but conclusion $\neg S(1,1)$ is False.
\end{tcolorbox}

\begin{tcolorbox}[title=(c) Universal-Existential Shift I]
\textbf{Valid.} Assume $\forall x(P(x) \to \exists y Q(y))$. For an arbitrary $x_0$, if $P(x_0)$ is false, the implication $P(x_0) \to Q(y)$ is vacuously true for any $y$. If $P(x_0)$ is true, there exists $y$ such that $Q(y)$ is true, making the implication true.
\end{tcolorbox}

\begin{tcolorbox}[title=(d) Universal-Existential Shift II]
\textbf{Valid.} Assume $\forall x \exists y(P(x) \to Q(y))$. For any $x$, there is a $y$ where $P(x) \to Q(y)$ holds. If $P(x)$ is true, $Q(y)$ must be true for that $y$. Thus $\exists y Q(y)$ is satisfied.
\end{tcolorbox}

\begin{tcolorbox}[title=(f) Antisymmetry vs. Irreflexivity]
\textbf{Invalid.} Let $D = \{0\}, S = \{(0,0)\}$. Premise $(S(0,0) \to 0=0)$ is true, but conclusion $\neg S(0,0)$ is false.
\end{tcolorbox}

\begin{tcolorbox}[title=(h) Global Equivalence]
\textbf{Invalid.} Let $D = \{0, 1\}, P = \{0\}$. For $x=0, y=1$, the conjunct $(P(0) \to P(1))$ is $T \to F$, which is false.
\end{tcolorbox}

\begin{tcolorbox}[title=(i) Bi-implication Distribution]
\textbf{Valid.} Assume $\forall x (P(x) \leftrightarrow Q(x))$ and $\forall x P(x)$. For any $a$, $P(a)$ is true. Since $P(a) \leftrightarrow Q(a)$, then $Q(a)$ must be true. Therefore, $\forall x Q(x)$ holds.
\end{tcolorbox}

\end{document}